\documentclass[11pt,a4paper]{article}
\usepackage[utf8]{inputenc}
\usepackage[T1]{fontenc}
\usepackage[catalan]{babel}
\usepackage{lmodern}
\usepackage[sfdefault,lining]{FiraSans}
\usepackage[a4paper,top=2.3cm,bottom=2.2cm,left=1.8cm,right=1.8cm,headheight=1.3cm,headsep=0.4cm,footskip=1cm]{geometry}
\usepackage{amsmath,amssymb}
\usepackage{icomma}
\usepackage{graphicx}
\usepackage[table]{xcolor}
\usepackage{tikz}
\usepackage{fontawesome5}
\usepackage{ccicons}
\usepackage{fancyhdr}
\usepackage{lastpage}
\usepackage{tasks}
\usepackage{needspace}
\usepackage[most]{tcolorbox}
\usepackage{array}
\usepackage[hidelinks,pdftitle={Solucions de la fitxa de la unitat 10: Estadística. 1r de Batxillerat MI},pdfsubject={Matemàtiques I},pdfkeywords={estadística, freqüències, mitjana, mediana, quartils, variància, desviació típica, coeficient de variació, variable bidimensional, taula de doble entrada, núvol de punts, covariància, correlació lineal}]{hyperref}

\definecolor{navy}{HTML}{1B2748}
\definecolor{accent}{HTML}{E8603A}
\definecolor{teal}{HTML}{1F9E90}
\definecolor{soft}{HTML}{EEF1F8}
\definecolor{muted}{HTML}{5D6685}

\setlength{\parindent}{0pt}
\setlength{\parskip}{3pt}
\settasks{label=\alph*), label-format=\bfseries\color{navy}, label-width=1.3em, item-indent=1.8em, column-sep=0.6em, after-item-skip=5pt, before-skip=3pt, after-skip=4pt}

% ---------- capçalera i peu ----------
\newcommand{\llicencia}{\href{https://creativecommons.org/licenses/by-nc-sa/4.0/deed.ca}{\ccbyncsa}\ \ Llicència Creative Commons Reconeixement-NoComercial-CompartirIgual 4.0 Internacional (CC BY-NC-SA 4.0)}
\pagestyle{fancy}
\fancyhf{}
\fancyhead[L]{\raisebox{-0.35\height}{\includegraphics[height=0.95cm]{logo-gva.png}}}
\fancyhead[C]{\small\color{muted}Solucions de la fitxa · Unitat 10: Estadística · 1r Batxillerat MI}
\fancyhead[R]{\raisebox{-0.35\height}{\includegraphics[height=0.95cm]{logo-vermellar.png}}}
\fancyfoot[L]{\scriptsize\color{muted}\llicencia}
\fancyfoot[R]{\small\color{muted}\thepage}
\renewcommand{\headrulewidth}{0.4pt}
\renewcommand{\headrule}{\color{muted!40}\hrule height \headrulewidth}
\fancypagestyle{primera}{\fancyhf{}\renewcommand{\headrulewidth}{0pt}
  \fancyfoot[L]{\scriptsize\color{muted}\llicencia}
  \fancyfoot[R]{\small\color{muted}\thepage}}

% ---------- nivells i exercicis ----------
\newcommand{\estrelles}[1]{{\color{accent}\ifcase#1\or\faStar\faStar[regular]\faStar[regular]\or\faStar\faStar\faStar[regular]\or\faStar\faStar\faStar\fi}}
\newcounter{exer}
\newcommand{\ex}[1]{\par\needspace{4\baselineskip}\medskip\refstepcounter{exer}\noindent{\color{navy}\bfseries\theexer.}\ \ }

\newcommand{\seccio}[2]{\par\needspace{16\baselineskip}\bigskip
  \begin{tcolorbox}[colback=navy,colframe=navy,coltext=white,arc=2mm,boxrule=0pt,left=3mm,right=3mm,top=1.2mm,bottom=1.2mm]
  {\large\bfseries #1\quad #2}
  \end{tcolorbox}\vspace{-2pt}}

\newcommand{\exf}[1]{\par\medskip\refstepcounter{exer}\noindent{\color{navy}\bfseries\theexer.}\ \ }
\providecommand{\gt}{>}\providecommand{\lt}{<}
% quadre de recordatori (només a la fitxa d'exercicis)
\newcommand{\recorda}[1]{\par\smallskip\begin{tcolorbox}[colback=accent!7,colframe=accent!55,arc=2mm,boxrule=0.5pt,left=2.5mm,right=2.5mm,top=1mm,bottom=1mm]\small{\color{accent}\faLightbulb[regular]}\ \textbf{Recorda:}\ #1\end{tcolorbox}}
\newcommand{\res}[1]{\boxed{#1}}\newcommand{\rest}[1]{\fbox{#1}}

\def\fmneed{11}% que el títol d'una secció no quede sol al final del full
% =====================================================================
% Capçalera FentMates de les fitxes: un color per a cada curs.
% Els colors són els de la pàgina inicial (LEVELS a assets/js/materials.js);
% el logo ve de _eines/logo/pdf/ (python3 _eines/logo/genera.py).
%
% Ús, al final del preàmbul de la unitat:
%   % =====================================================================
% Capçalera FentMates de les fitxes: un color per a cada curs.
% Els colors són els de la pàgina inicial (LEVELS a assets/js/materials.js);
% el logo ve de _eines/logo/pdf/ (python3 _eines/logo/genera.py).
%
% Ús, al final del preàmbul de la unitat:
%   % =====================================================================
% Capçalera FentMates de les fitxes: un color per a cada curs.
% Els colors són els de la pàgina inicial (LEVELS a assets/js/materials.js);
% el logo ve de _eines/logo/pdf/ (python3 _eines/logo/genera.py).
%
% Ús, al final del preàmbul de la unitat:
%   \input{../comu/fentmates.tex}
%   \fmcurs{1eso}
% i, al començament del document:
%   \capcalera{Matemàtiques}{1r ESO}{Fitxa d'exercicis · Unitat 2}{Divisibilitat}
% Opcional abans de \input: \def\fmneed{14} (espai mínim abans de cada secció).
% =====================================================================
\makeatletter
\def\fm@logodir{../../logo/pdf}
\providecommand{\fmneed}{10}
\def\fm@cursnom{}

% clau del curs, color i color del text damunt del color
\newcommand{\fm@curs}[3]{\@namedef{fm@color@#1}{#2}\@namedef{fm@text@#1}{#3}}
\fm@curs{1eso}{E6B84A}{navy}
\fm@curs{2eso}{3B82F6}{white}
\fm@curs{3eso}{FF7A50}{white}
\fm@curs{4eso-a}{A78BFA}{white}
\fm@curs{4eso-b}{6D28D9}{white}
\fm@curs{1batx-mi}{9B1C1C}{white}
\fm@curs{1batx-ccss}{EF6B6B}{white}
\fm@curs{2batx-mii}{1E6B47}{white}
\fm@curs{2batx-ccss}{2BB3A3}{white}

\newcommand{\fmcurs}[1]{%
  \@ifundefined{fm@color@#1}{\PackageError{fentmates}{Curs desconegut: #1}{}}{}%
  \edef\fm@tmp{\noexpand\definecolor{curs}{HTML}{\@nameuse{fm@color@#1}}%
               \noexpand\colorlet{curstext}{\@nameuse{fm@text@#1}}}\fm@tmp
  \colorlet{cursfosc}{curs!72!black}%
  \def\fm@logo{\fm@logodir/vertical-#1.pdf}%
  \def\fm@marca{\fm@logodir/marca-#1.pdf}}

% ---------- capçalera de la primera pàgina ----------
% #1 matèria (petit), #2 curs, #3 tipus de document i unitat, #4 títol
\newcommand{\capcalera}[4]{%
  \gdef\fm@cursnom{#2}%
  \par\noindent
  \begin{tcolorbox}[colback=curs,colframe=curs,boxrule=0pt,arc=2.2mm,width=3.5cm,height=2.15cm,
      valign=center,left=3mm,right=2mm,top=1mm,bottom=1mm,nobeforeafter,box align=center]
    \color{curstext}{\footnotesize #1}\par\vspace{1pt}{\fontsize{19}{21}\selectfont\bfseries #2\par}%
  \end{tcolorbox}\hspace{5mm}%
  \begin{minipage}[c]{\dimexpr\linewidth-3.5cm-5mm-2.9cm\relax}
    {\color{cursfosc}\small\bfseries\MakeUppercase{#3}}\par\vspace{2pt}
    {\color{navy}\fontsize{20}{23}\selectfont\bfseries #4\par}
  \end{minipage}\hfill
  \begin{minipage}[c]{2.6cm}\raggedleft\includegraphics[width=2.6cm]{\fm@logo}\end{minipage}\par
  \vspace{2pt}}

% ---------- barres de secció del color del curs ----------
\renewcommand{\seccio}[2]{\par\needspace{\fmneed\baselineskip}\bigskip
  \begin{tcolorbox}[colback=curs,colframe=curs,coltext=curstext,arc=2mm,boxrule=0pt,left=3mm,right=3mm,top=1.2mm,bottom=1.2mm]
  {\large\bfseries #1\quad #2}
  \end{tcolorbox}\vspace{-2pt}}

% ---------- estrelles de dificultat del color del curs (a Batxillerat no se'n posen) ----------
\@ifundefined{estrelles}{}{%
  \renewcommand{\estrelles}[1]{{\color{curs}\ifcase#1\or\faStar\faStar[regular]\faStar[regular]\or\faStar\faStar\faStar[regular]\or\faStar\faStar\faStar\fi}}}

% ---------- capçalera i peu de les pàgines ----------
% primera pàgina: logos institucionals al peu
\fancypagestyle{primera}{\fancyhf{}\renewcommand{\headrulewidth}{0pt}
  \fancyfoot[L]{\raisebox{-0.32\height}{\includegraphics[height=0.78cm]{logo-gva.png}}\hspace{5mm}%
    \raisebox{-0.32\height}{\includegraphics[height=0.78cm]{logo-vermellar.png}}}
  \fancyfoot[R]{\scriptsize\color{muted}\href{https://creativecommons.org/licenses/by-nc-sa/4.0/deed.ca}{\ccbyncsa}\ \ CC BY-NC-SA 4.0 · fentmates.org\qquad\small\thepage}}
% la resta: el símbol a l'esquerra i el curs a la dreta, amb el títol al centre com abans
\fancyhead[L]{\raisebox{-0.3\height}{\includegraphics[height=0.62cm]{\fm@marca}}}
\fancyhead[R]{\small\bfseries\color{cursfosc}\fm@cursnom}
\renewcommand{\headrule}{\color{curs!45}\hrule height \headrulewidth}
\makeatother

%   \fmcurs{1eso}
% i, al començament del document:
%   \capcalera{Matemàtiques}{1r ESO}{Fitxa d'exercicis · Unitat 2}{Divisibilitat}
% Opcional abans de \input: \def\fmneed{14} (espai mínim abans de cada secció).
% =====================================================================
\makeatletter
\def\fm@logodir{../../logo/pdf}
\providecommand{\fmneed}{10}
\def\fm@cursnom{}

% clau del curs, color i color del text damunt del color
\newcommand{\fm@curs}[3]{\@namedef{fm@color@#1}{#2}\@namedef{fm@text@#1}{#3}}
\fm@curs{1eso}{E6B84A}{navy}
\fm@curs{2eso}{3B82F6}{white}
\fm@curs{3eso}{FF7A50}{white}
\fm@curs{4eso-a}{A78BFA}{white}
\fm@curs{4eso-b}{6D28D9}{white}
\fm@curs{1batx-mi}{9B1C1C}{white}
\fm@curs{1batx-ccss}{EF6B6B}{white}
\fm@curs{2batx-mii}{1E6B47}{white}
\fm@curs{2batx-ccss}{2BB3A3}{white}

\newcommand{\fmcurs}[1]{%
  \@ifundefined{fm@color@#1}{\PackageError{fentmates}{Curs desconegut: #1}{}}{}%
  \edef\fm@tmp{\noexpand\definecolor{curs}{HTML}{\@nameuse{fm@color@#1}}%
               \noexpand\colorlet{curstext}{\@nameuse{fm@text@#1}}}\fm@tmp
  \colorlet{cursfosc}{curs!72!black}%
  \def\fm@logo{\fm@logodir/vertical-#1.pdf}%
  \def\fm@marca{\fm@logodir/marca-#1.pdf}}

% ---------- capçalera de la primera pàgina ----------
% #1 matèria (petit), #2 curs, #3 tipus de document i unitat, #4 títol
\newcommand{\capcalera}[4]{%
  \gdef\fm@cursnom{#2}%
  \par\noindent
  \begin{tcolorbox}[colback=curs,colframe=curs,boxrule=0pt,arc=2.2mm,width=3.5cm,height=2.15cm,
      valign=center,left=3mm,right=2mm,top=1mm,bottom=1mm,nobeforeafter,box align=center]
    \color{curstext}{\footnotesize #1}\par\vspace{1pt}{\fontsize{19}{21}\selectfont\bfseries #2\par}%
  \end{tcolorbox}\hspace{5mm}%
  \begin{minipage}[c]{\dimexpr\linewidth-3.5cm-5mm-2.9cm\relax}
    {\color{cursfosc}\small\bfseries\MakeUppercase{#3}}\par\vspace{2pt}
    {\color{navy}\fontsize{20}{23}\selectfont\bfseries #4\par}
  \end{minipage}\hfill
  \begin{minipage}[c]{2.6cm}\raggedleft\includegraphics[width=2.6cm]{\fm@logo}\end{minipage}\par
  \vspace{2pt}}

% ---------- barres de secció del color del curs ----------
\renewcommand{\seccio}[2]{\par\needspace{\fmneed\baselineskip}\bigskip
  \begin{tcolorbox}[colback=curs,colframe=curs,coltext=curstext,arc=2mm,boxrule=0pt,left=3mm,right=3mm,top=1.2mm,bottom=1.2mm]
  {\large\bfseries #1\quad #2}
  \end{tcolorbox}\vspace{-2pt}}

% ---------- estrelles de dificultat del color del curs (a Batxillerat no se'n posen) ----------
\@ifundefined{estrelles}{}{%
  \renewcommand{\estrelles}[1]{{\color{curs}\ifcase#1\or\faStar\faStar[regular]\faStar[regular]\or\faStar\faStar\faStar[regular]\or\faStar\faStar\faStar\fi}}}

% ---------- capçalera i peu de les pàgines ----------
% primera pàgina: logos institucionals al peu
\fancypagestyle{primera}{\fancyhf{}\renewcommand{\headrulewidth}{0pt}
  \fancyfoot[L]{\raisebox{-0.32\height}{\includegraphics[height=0.78cm]{logo-gva.png}}\hspace{5mm}%
    \raisebox{-0.32\height}{\includegraphics[height=0.78cm]{logo-vermellar.png}}}
  \fancyfoot[R]{\scriptsize\color{muted}\href{https://creativecommons.org/licenses/by-nc-sa/4.0/deed.ca}{\ccbyncsa}\ \ CC BY-NC-SA 4.0 · fentmates.org\qquad\small\thepage}}
% la resta: el símbol a l'esquerra i el curs a la dreta, amb el títol al centre com abans
\fancyhead[L]{\raisebox{-0.3\height}{\includegraphics[height=0.62cm]{\fm@marca}}}
\fancyhead[R]{\small\bfseries\color{cursfosc}\fm@cursnom}
\renewcommand{\headrule}{\color{curs!45}\hrule height \headrulewidth}
\makeatother

%   \fmcurs{1eso}
% i, al començament del document:
%   \capcalera{Matemàtiques}{1r ESO}{Fitxa d'exercicis · Unitat 2}{Divisibilitat}
% Opcional abans de \input: \def\fmneed{14} (espai mínim abans de cada secció).
% =====================================================================
\makeatletter
\def\fm@logodir{../../logo/pdf}
\providecommand{\fmneed}{10}
\def\fm@cursnom{}

% clau del curs, color i color del text damunt del color
\newcommand{\fm@curs}[3]{\@namedef{fm@color@#1}{#2}\@namedef{fm@text@#1}{#3}}
\fm@curs{1eso}{E6B84A}{navy}
\fm@curs{2eso}{3B82F6}{white}
\fm@curs{3eso}{FF7A50}{white}
\fm@curs{4eso-a}{A78BFA}{white}
\fm@curs{4eso-b}{6D28D9}{white}
\fm@curs{1batx-mi}{9B1C1C}{white}
\fm@curs{1batx-ccss}{EF6B6B}{white}
\fm@curs{2batx-mii}{1E6B47}{white}
\fm@curs{2batx-ccss}{2BB3A3}{white}

\newcommand{\fmcurs}[1]{%
  \@ifundefined{fm@color@#1}{\PackageError{fentmates}{Curs desconegut: #1}{}}{}%
  \edef\fm@tmp{\noexpand\definecolor{curs}{HTML}{\@nameuse{fm@color@#1}}%
               \noexpand\colorlet{curstext}{\@nameuse{fm@text@#1}}}\fm@tmp
  \colorlet{cursfosc}{curs!72!black}%
  \def\fm@logo{\fm@logodir/vertical-#1.pdf}%
  \def\fm@marca{\fm@logodir/marca-#1.pdf}}

% ---------- capçalera de la primera pàgina ----------
% #1 matèria (petit), #2 curs, #3 tipus de document i unitat, #4 títol
\newcommand{\capcalera}[4]{%
  \gdef\fm@cursnom{#2}%
  \par\noindent
  \begin{tcolorbox}[colback=curs,colframe=curs,boxrule=0pt,arc=2.2mm,width=3.5cm,height=2.15cm,
      valign=center,left=3mm,right=2mm,top=1mm,bottom=1mm,nobeforeafter,box align=center]
    \color{curstext}{\footnotesize #1}\par\vspace{1pt}{\fontsize{19}{21}\selectfont\bfseries #2\par}%
  \end{tcolorbox}\hspace{5mm}%
  \begin{minipage}[c]{\dimexpr\linewidth-3.5cm-5mm-2.9cm\relax}
    {\color{cursfosc}\small\bfseries\MakeUppercase{#3}}\par\vspace{2pt}
    {\color{navy}\fontsize{20}{23}\selectfont\bfseries #4\par}
  \end{minipage}\hfill
  \begin{minipage}[c]{2.6cm}\raggedleft\includegraphics[width=2.6cm]{\fm@logo}\end{minipage}\par
  \vspace{2pt}}

% ---------- barres de secció del color del curs ----------
\renewcommand{\seccio}[2]{\par\needspace{\fmneed\baselineskip}\bigskip
  \begin{tcolorbox}[colback=curs,colframe=curs,coltext=curstext,arc=2mm,boxrule=0pt,left=3mm,right=3mm,top=1.2mm,bottom=1.2mm]
  {\large\bfseries #1\quad #2}
  \end{tcolorbox}\vspace{-2pt}}

% ---------- estrelles de dificultat del color del curs (a Batxillerat no se'n posen) ----------
\@ifundefined{estrelles}{}{%
  \renewcommand{\estrelles}[1]{{\color{curs}\ifcase#1\or\faStar\faStar[regular]\faStar[regular]\or\faStar\faStar\faStar[regular]\or\faStar\faStar\faStar\fi}}}

% ---------- capçalera i peu de les pàgines ----------
% primera pàgina: logos institucionals al peu
\fancypagestyle{primera}{\fancyhf{}\renewcommand{\headrulewidth}{0pt}
  \fancyfoot[L]{\raisebox{-0.32\height}{\includegraphics[height=0.78cm]{logo-gva.png}}\hspace{5mm}%
    \raisebox{-0.32\height}{\includegraphics[height=0.78cm]{logo-vermellar.png}}}
  \fancyfoot[R]{\scriptsize\color{muted}\href{https://creativecommons.org/licenses/by-nc-sa/4.0/deed.ca}{\ccbyncsa}\ \ CC BY-NC-SA 4.0 · fentmates.org\qquad\small\thepage}}
% la resta: el símbol a l'esquerra i el curs a la dreta, amb el títol al centre com abans
\fancyhead[L]{\raisebox{-0.3\height}{\includegraphics[height=0.62cm]{\fm@marca}}}
\fancyhead[R]{\small\bfseries\color{cursfosc}\fm@cursnom}
\renewcommand{\headrule}{\color{curs!45}\hrule height \headrulewidth}
\makeatother

\fmcurs{1batx-mi}

\begin{document}
\thispagestyle{primera}

\capcalera{Matemàtiques I}{1r Batx}{Solucions de la fitxa · Unitat 10}{\raggedright Estadística}
\vspace{8pt}

\begin{tcolorbox}[colback=soft,colframe=soft,arc=2mm,left=3mm,right=3mm,top=1.5mm,bottom=1.5mm]
\small Intenta resoldre cada exercici abans de mirar-ne la solució. Si no t'ix el mateix resultat, repassa els passos o comprova'l amb les calculadores dels apunts interactius de la unitat.
\end{tcolorbox}

% =====================================================
\seccio{1}{Estadística unidimensional: variables i taules}

\exf{1} \begin{tasks}(2)
  \task \rest{qualitativa}
  \task \rest{quantitativa discreta}
  \task \rest{quantitativa contínua}
  \task \rest{quantitativa contínua}
  \task \rest{qualitativa}
  \task \rest{quantitativa discreta}
\end{tasks}

\ex{1} \par{\centering {\small$\begin{array}{|c|c|c|c|c|}\hline x_i & f_i & F_i & h_i & H_i\\\hline 3 & 1 & 1 & 0{,}04 & 0{,}04\\ 4 & 2 & 3 & 0{,}08 & 0{,}12\\ 5 & 5 & 8 & 0{,}2 & 0{,}32\\ 6 & 8 & 16 & 0{,}32 & 0{,}64\\ 7 & 5 & 21 & 0{,}2 & 0{,}84\\ 8 & 3 & 24 & 0{,}12 & 0{,}96\\ 9 & 1 & 25 & 0{,}04 & 1\\\hline\end{array}$}\par}\smallskip Han aprovat $25-3=22$ alumnes: el $\res{88\,\%}$.


% =====================================================
\seccio{2}{Paràmetres de centralització i de posició}

\exf{1} \begin{tasks}(1)
  \task $\bar x=\frac{51}{9}\approx\res{5{,}67}$, \ $Me=\res{6}$, \ $Mo=\res{7}$, \ $Q_1=\res{4}$, \ $Q_3=\res{7}$
  \task $\bar x=\frac{156}{8}=\res{19{,}5}$, \ $Me=\frac{18+20}{2}=\res{19}$, \ $Mo=\res{15}$, \ $Q_1=\res{15}$, \ $Q_3=\res{23}$
\end{tasks}

\ex{1} $\bar x=\frac{152}{25}=\res{6{,}08}$; \ $\frac N2=12{,}5$: $Me=\res{6}$; \ $Mo=\res{6}$; \ $Q_1=\res{5}$ ($\frac N4=6{,}25$), \ $Q_3=\res{7}$ ($\frac{3N}4=18{,}75$).

\ex{2} Ara suma $4\cdot6{,}5=26$; ha de sumar $5\cdot7=35$: necessita un $\res{9}$.


% =====================================================
\seccio{3}{Paràmetres de dispersió}

\exf{1} \begin{tasks}(1)
  \task A: $\bar x=\res{8}$, rang $\res{4}$, $\sigma^2=\frac{588}{9}-64\approx\res{1{,}33}$, $\sigma\approx\res{1{,}15}$
  \task B: $\bar x=\res{8}$, rang $\res{12}$, $\sigma^2=\frac{688}{9}-64\approx\res{12{,}44}$, $\sigma\approx\res{3{,}53}$
\end{tasks}\par Tenen la mateixa mitjana, però les notes del grup B estan \rest{molt més disperses}.

\ex{2} Marques de classe $5,\ 15,\ 25,\ 35,\ 45$. \ $\sum x_if_i=890$, \ $\sum x_i^2f_i=24\,400$. \ $\bar x=\frac{890}{40}=\res{22{,}25}$ min, \ $\sigma^2\approx114{,}94$, \ $\sigma\approx\res{10{,}72}$ min, \ $CV\approx\res{0{,}482}$ ($48{,}2\,\%$).

\ex{2} \begin{tasks}(1)
  \task A: $\bar x=12{,}2$ s, \ $\sigma\approx0{,}141$ s
  \task B: $\bar x=12{,}2$ s, \ $\sigma\approx0{,}469$ s
\end{tasks}\par Tenen la mateixa mitjana; l'atleta \rest{A és més regular} (menys dispersió).

\ex{2} Tenen unitats i mitjanes diferents: cal el coeficient de variació. \ $CV_1=\frac{8}{165}\approx0{,}048$, \ $CV_2=\frac{4}{32}=0{,}125$: és més dispers \rest{el segon grup} (pesos).


% =====================================================
\seccio{4}{Variables bidimensionals. Taules de doble entrada}

\exf{1} \par{\centering {\small$\begin{array}{|c|c|c|c|c|}\hline X\backslash Y & 0 & 1 & 2 & f_x\\\hline 1 & 4 & 2 & 0 & 6\\ 2 & 3 & 6 & 3 & 12\\ 3 & 0 & 3 & 4 & 7\\\hline f_y & 7 & 11 & 7 & 25\\\hline \end{array}$}\par}\smallskip\begin{tasks}(1)
  \task Vegeu la taula.
  \task $\res{6}$ famílies amb 2 fills i 1 cotxe; \ $\res{7}$ amb més d'un cotxe (2 cotxes).
  \task $\bar x=\frac{51}{25}=\res{2{,}04}$, \ $\bar y=\frac{25}{25}=\res{1}$
\end{tasks}

\ex{2} \par{\centering {\small$\begin{array}{|c|c|c|c|c|}\hline X\backslash Y & 1 & 2 & 3 & f_x\\\hline 1 & 2 & 1 & 0 & 3\\ 2 & 2 & 3 & 0 & 5\\ 3 & 0 & 2 & 2 & 4\\\hline f_y & 4 & 6 & 2 & 12\\\hline \end{array}$}\par}\smallskip $\bar x=\frac{25}{12}\approx\res{2{,}08}$, \ $\bar y=\frac{22}{12}\approx\res{1{,}83}$.


% =====================================================
\seccio{5}{Núvol de punts}

\exf{1} $(\bar x,\bar y)=(4{,}5,\ 51{,}25)$. \ Els punts baixen gairebé en línia recta: correlació \rest{negativa i molt forta}.\par\smallskip{\centering \begin{tikzpicture}[scale=0.42,baseline=(current bounding box.center)]\draw[muted!30,very thin] (0,0) grid (9.0,6.0);\draw[-latex,thick] (0,0)--(9.5,0) node[right,font=\scriptsize]{$h$};\draw[-latex,thick] (0,0)--(0,6.5) node[above,font=\scriptsize]{$l$};\node[below,font=\tiny] at (0.0,0) {$0$};\node[below,font=\tiny] at (1.0,0) {$1$};\node[below,font=\tiny] at (2.0,0) {$2$};\node[below,font=\tiny] at (3.0,0) {$3$};\node[below,font=\tiny] at (4.0,0) {$4$};\node[below,font=\tiny] at (5.0,0) {$5$};\node[below,font=\tiny] at (6.0,0) {$6$};\node[below,font=\tiny] at (7.0,0) {$7$};\node[below,font=\tiny] at (8.0,0) {$8$};\node[below,font=\tiny] at (9.0,0) {$9$};\node[left,font=\tiny] at (0,0.0) {$35$};\node[left,font=\tiny] at (0,1.0) {$40$};\node[left,font=\tiny] at (0,2.0) {$45$};\node[left,font=\tiny] at (0,3.0) {$50$};\node[left,font=\tiny] at (0,4.0) {$55$};\node[left,font=\tiny] at (0,5.0) {$60$};\node[left,font=\tiny] at (0,6.0) {$65$};\fill[accent] (1.0,5.4) circle (0.16);\fill[accent] (2.0,5.0) circle (0.16);\fill[accent] (3.0,4.0) circle (0.16);\fill[accent] (4.0,3.6) circle (0.16);\fill[accent] (5.0,3.0) circle (0.16);\fill[accent] (6.0,2.2) circle (0.16);\fill[accent] (7.0,1.8) circle (0.16);\fill[accent] (8.0,1.0) circle (0.16);\draw[navy,thick,fill=white] (9/2,13/4) circle (0.2);\end{tikzpicture}\par}

\ex{1} \begin{tasks}(4)
  \task $\res{r=0{,}9}$
  \task $\res{r=-0{,}95}$
  \task $\res{r=0}$
  \task $\res{r=-0{,}5}$
\end{tasks}

\ex{1} \begin{tasks}(2)
  \task \rest{positiva}
  \task \rest{negativa}
  \task \rest{cap}
  \task \rest{positiva}
  \task \rest{negativa}
  \task \rest{cap}
\end{tasks}


% =====================================================
\seccio{6}{Variàncies i covariància}

\exf{2} \par{\centering {\small$\begin{array}{|c|c|c|c|c|}\hline x_i & y_i & x_i^2 & y_i^2 & x_iy_i\\\hline 2 & 3 & 4 & 9 & 6\\ 4 & 5 & 16 & 25 & 20\\ 5 & 4 & 25 & 16 & 20\\ 6 & 7 & 36 & 49 & 42\\ 8 & 9 & 64 & 81 & 72\\\hline \textstyle\sum: 25 & 28 & 145 & 180 & 160\\\hline\end{array}$}\par}\smallskip $\bar x=\frac{25}{5}=5$, \ $\bar y=\frac{28}{5}=5{,}6$; \ $\sigma_x^2=\frac{145}{5}-\bar x^2=4$, \ $\sigma_y^2=\frac{180}{5}-\bar y^2=4{,}64$; \ $\sigma_x=2$, \ $\sigma_y\approx2{,}154$. \ $\sigma_{xy}=\frac{160}{5}-\bar x\,\bar y=\res{4}$.

\ex{2} $\bar x=\frac{150}{5}=30$, \ $\bar y=\frac{30}{5}=6$; \ $\sigma_x^2=\frac{5500}{5}-\bar x^2=200$, \ $\sigma_y^2=\frac{210}{5}-\bar y^2=6$; \ $\sigma_x\approx14{,}142$, \ $\sigma_y\approx2{,}449$. \ $\sigma_{xy}=\frac{730}{5}-\bar x\,\bar y=\res{-34}$. \ És negativa: quan $X$ augmenta, $Y$ tendeix a \rest{disminuir}.

\ex{2} $\sum x_iy_jf_{ij}=50$: \ $\sigma_{xy}=\frac{50}{12}-\bar x\,\bar y\approx\res{0{,}347}$

\ex{2} $\sum x_iy_jf_{ij}=59$: \ $\sigma_{xy}=\frac{59}{25}-2{,}04\cdot1=\res{0{,}32}$


% =====================================================
\seccio{7}{Correlació lineal}

\exf{2} \begin{tasks}(1)
  \task $r=\dfrac{\sigma_{xy}}{\sigma_x\sigma_y}\approx\res{0{,}928}$: correlació \rest{positiva forta}.
  \task $r=\dfrac{\sigma_{xy}}{\sigma_x\sigma_y}\approx\res{-0{,}981}$: correlació \rest{negativa molt forta}.
  \task $r=\dfrac{0{,}32}{\sqrt{0{,}518}\cdot\sqrt{0{,}56}}\approx\res{0{,}594}$: correlació \rest{positiva moderada}.
\end{tasks}

\ex{2} $\bar x=\frac{1155}{7}=165$, \ $\bar y=\frac{421}{7}\approx60{,}143$; \ $\sigma_x^2=\frac{191\,275}{7}-\bar x^2=100$, \ $\sigma_y^2=\frac{25\,807}{7}-\bar y^2\approx69{,}551$; \ $\sigma_x=10$, \ $\sigma_y\approx8{,}34$. \ $\sigma_{xy}=\frac{70\,045}{7}-\bar x\,\bar y\approx\res{82{,}857}$. $r=\dfrac{\sigma_{xy}}{\sigma_x\sigma_y}\approx\res{0{,}994}$: correlació \rest{positiva molt forta}.\par\smallskip{\centering \begin{tikzpicture}[scale=0.42,baseline=(current bounding box.center)]\draw[muted!30,very thin] (0,0) grid (8.0,6.0);\draw[-latex,thick] (0,0)--(8.5,0) node[right,font=\scriptsize]{$cm$};\draw[-latex,thick] (0,0)--(0,6.5) node[above,font=\scriptsize]{$kg$};\node[below,font=\tiny] at (0.0,0) {$145$};\node[below,font=\tiny] at (2.0,0) {$155$};\node[below,font=\tiny] at (4.0,0) {$165$};\node[below,font=\tiny] at (6.0,0) {$175$};\node[below,font=\tiny] at (8.0,0) {$185$};\node[left,font=\tiny] at (0,0.0) {$45$};\node[left,font=\tiny] at (0,1.0) {$50$};\node[left,font=\tiny] at (0,2.0) {$55$};\node[left,font=\tiny] at (0,3.0) {$60$};\node[left,font=\tiny] at (0,4.0) {$65$};\node[left,font=\tiny] at (0,5.0) {$70$};\node[left,font=\tiny] at (0,6.0) {$75$};\fill[accent] (1.0,0.6) circle (0.16);\fill[accent] (2.0,1.4) circle (0.16);\fill[accent] (3.0,2.0) circle (0.16);\fill[accent] (4.0,3.2) circle (0.16);\fill[accent] (5.0,3.6) circle (0.16);\fill[accent] (6.0,5.0) circle (0.16);\fill[accent] (7.0,5.4) circle (0.16);\end{tikzpicture}\par}

\ex{2} $\bar x=\frac{15}{6}=2{,}5$, \ $\bar y=\frac{40}{6}\approx6{,}667$; \ $\sigma_x^2=\frac{55}{6}-\bar x^2\approx2{,}917$, \ $\sigma_y^2=\frac{306}{6}-\bar y^2\approx6{,}556$; \ $\sigma_x\approx1{,}708$, \ $\sigma_y\approx2{,}56$. \ $\sigma_{xy}=\frac{75}{6}-\bar x\,\bar y\approx\res{-4{,}167}$. $r=\dfrac{\sigma_{xy}}{\sigma_x\sigma_y}\approx\res{-0{,}953}$: correlació \rest{negativa molt forta}. \ Però una correlació, encara que siga forta, \rest{no demostra que hi haja una relació de causa i efecte}: hi pot haver altres variables (el temps d'estudi, per exemple).

\ex{2} $\bar x=\frac{153}{6}=25{,}5$, \ $\bar y=\frac{406}{6}\approx67{,}667$; \ $\sigma_x^2=\frac{4059}{6}-\bar x^2=26{,}25$, \ $\sigma_y^2=\frac{30\,378}{6}-\bar y^2\approx484{,}222$; \ $\sigma_x\approx5{,}123$, \ $\sigma_y\approx22{,}005$. \ $\sigma_{xy}=\frac{11\,025}{6}-\bar x\,\bar y=\res{112}$. $r=\dfrac{\sigma_{xy}}{\sigma_x\sigma_y}\approx\res{0{,}993}$: correlació \rest{positiva molt forta}.

\ex{2} $\sigma_x^2\approx0{,}576$, \ $\sigma_y^2\approx0{,}472$, \ $\sigma_{xy}\approx0{,}347$. \ $r\approx\res{0{,}666}$: correlació \rest{positiva moderada}.

\ex{1} $\sigma_{xy}=r\,\sigma_x\,\sigma_y=0{,}8\cdot2\cdot5=\res{8}$.

\ex{1} $r=\frac{-12}{4\cdot5}=\res{-0{,}6}$ (negativa moderada). \ En el segon cas eixiria $r=\frac{3}{2}=1{,}5>1$: \rest{no és possible}, hi ha un error en les dades.

\ex{3} $\bar x=\frac{50}{10}=5$, \ $\bar y=\frac{80}{10}=8$; \ $\sigma_x^2=\frac{290}{10}-\bar x^2=4$, \ $\sigma_y^2=\frac{730}{10}-\bar y^2=9$; \ $\sigma_x=2$, \ $\sigma_y=3$. \ $\sigma_{xy}=\frac{440}{10}-\bar x\,\bar y=\res{4}$. $r=\dfrac{\sigma_{xy}}{\sigma_x\sigma_y}\approx\res{0{,}667}$: correlació \rest{positiva moderada}.


% =====================================================
\seccio{8}{Altres tipus de relació. Correlació i causalitat}

\exf{2} \begin{tasks}(1)
  \task Relació \rest{parabòlica} (creix i després decreix): és forta, però no és lineal i $r\approx0{,}02$, \rest{pròxim a 0}.
  \task Relació \rest{exponencial}: creixent, $r\approx0{,}9$ és alt, però els punts no segueixen una recta.
  \task Els punts de baix no tenen cap relació, però el \rest{valor atípic} $(9,9)$ fa que $r\approx0{,}99$. Cal revisar-lo abans de traure conclusions.
\end{tasks}

\ex{2} \rest{No}: la correlació no implica causalitat. Hi ha una tercera variable, la \rest{grandària de l'incendi}: com més gran és, més bombers hi van i més danys causa.

\ex{1} a) \rest{Falsa}: només vol dir que no hi ha relació lineal (pot ser parabòlica, per exemple). \ b) \rest{Falsa}: és molt forta; el signe indica que és inversa. \ c) \rest{Vertadera}.


\end{document}
